\documentclass[11pt]{article}

% ====================================================
% PACKAGES
% ====================================================
\usepackage[margin=1in]{geometry}
\usepackage[T1]{fontenc}
\usepackage{lmodern}
\usepackage{microtype}
\usepackage{setspace}
\usepackage{parskip}

\usepackage{amsmath, amssymb, mathtools}
\usepackage{siunitx}
\usepackage{graphicx}
\usepackage{float}
\usepackage[justification=centering]{caption}

\usepackage{enumitem}
\usepackage{booktabs}
\usepackage{tabularx}
\usepackage{multirow}
\usepackage{makecell}

\usepackage{xcolor}
\usepackage{hyperref}
\usepackage{xurl}

\usepackage{fancyhdr}
\usepackage{lastpage}

% ====================================================
% Assignment Info
% ====================================================
\newcommand{\Name}{Arael Anaya}
\newcommand{\Assignment}{Friction Characterization of the Scissor Boom}
\newcommand{\course}{Senior Design}

\newcommand{\dth}{\dot{\theta}}
\newcolumntype{Y}{>{\raggedright\arraybackslash}X}

% ====================================================
% PAGE STYLE
% ====================================================
\onehalfspacing%

\pagestyle{fancy}
\fancyhf{}
\lhead{\Assignment\ |\ \course}
\rhead{\Name\ |\ \today}
\cfoot{Page \thepage\ of~\pageref{LastPage}}

\hypersetup{
    colorlinks,
    linkcolor={red!50!black},
    citecolor={blue!50!black},
    urlcolor={blue!80!black}
}


% ====================================================
% TITLE
% ====================================================
\title{Friction Sources and Resistive-Torque Budget\\
\large Paragon Deployable Radiator --- Scissor Boom with Inflatable Kink Actuation}
\author{\Name\\ \course}
\date{\today}

\begin{document}
\maketitle

\section{Purpose}

The deployment model (\texttt{Derivation/main.tex}, \texttt{Phase1.py}, \texttt{Phase2.py}) represents
every resistive effect with one untested placeholder in \texttt{params.py}:
$F_{\text{friction}} = \SI{5}{N}$ per joint, applied at $r_o$ and multiplied by a
$1.5$ safety factor. This document replaces that guess with a component-by-component list of friction and
other parasitic resistances. Each value comes from a cited source. The list is summed into a system total
and converted back into the model's $F_{\text{friction}}$ convention so it can be dropped into the existing
scripts.

Every number in the tables below comes from \texttt{friction\_budget.py} (this folder). That script imports
the project's own \texttt{params.py}, \texttt{Phase1.py} and \texttt{friction\_sweep.py}. Run it once with
\texttt{USE\_TESTING\_PARAMS = False} (production, on orbit) and once with \texttt{True} (bench article, 1\,g).

\begin{figure}[H]
    \centering
    \includegraphics[width=0.5\linewidth]{Figures/systemDiagram.png}
    \caption{Scissor geometry: pinned at the root, root top node on a guide, $n$ bays of length $L$.}
\end{figure}

\section{Method}

\subsection*{Common currency: generalized torque about $\theta$}

The derivation's equation of motion carries friction as $Q_F = N_s\tau_f$, a generalized torque about
$\theta$. Every source below is converted to that same quantity $Q$ (N\,m) by virtual work. The sources can
then be summed and compared directly with the spring torque $N_s k\theta$.

\subsection*{Joint count and rotation}

One X per bay gives $3n$ revolute joints: $n$ centre pivots, $2(n-1)$ inter-bay nodes, the ground pin, and
the pin on the guide carriage. The centre pivots and inter-bay nodes join a $+\theta$ rod to a $-\theta$
rod, so they turn $2\,d\theta$. The ground and carriage pins turn $d\theta$. A per-joint friction torque
$\tau_j$ therefore costs
\begin{equation}
Q_{\text{pins}} = \tau_j\big[2(3n-2) + 2\big] = (6n-2)\,\tau_j ,
\qquad \tau_j = \mu N r_{\text{pin}}
\label{eq:pins}
\end{equation}
which is $40\,\tau_j$ for the production array ($n=7$) and $10\,\tau_j$ for the bench article ($n=2$). This
factor matters. The model's $N_s\tau_f$ counts only $n$ joints turning $d\theta$ each, so it under-counts
pin friction by $(6n-2)/n \approx 5.7\times$ for $n=7$.

\subsection*{Pin normal loads}

\begin{itemize}[nosep]
    \item \textbf{On orbit (production):} there is no weight, so pin loads come from the spring torque
    reacted across a half-link: $N \approx 2k\theta_0/L$. This is evaluated at stow, where it peaks.
    \item \textbf{Bench (1\,g):} every joint is assumed to carry one full bay's weight, $N = M_j g$. This is
    conservative and follows the 1\,g bushing-load approach of Guzik and Benafan~\cite{guzik}.
\end{itemize}

The spring rate $k$ is the value the current pipeline produces from the \SI{5}{N} assumption, using the
\texttt{required\_k} function in \texttt{friction\_sweep.py}: $k=\SI{0.397}{N.m/rad}$ for production and
$k=\SI{0.0248}{N.m/rad}$ for the bench. Friction terms that depend on $k$ (pins, guide and spring losses) are
small next to the harness term. Resizing $k$ therefore changes the nominal total by about 8\% for production
and 16\% for the bench (Section~\ref{sec:sc}).

\subsection*{Low / nominal / high}

Each source is given a range. \emph{Low} is the best case the literature supports. \emph{Nominal} is the
recommended design value for the intended hardware: self-lubricated bushings on the bench, MoS$_2$-lubricated
joints on orbit. \emph{High} is a credible worst case, such as bare metal-on-metal static friction or the
worst harness characterization. Each source is also tagged with its ECSS contributor class (FR friction,
HY hysteresis, HA harness, HD adhesion) so the standard uncertainty factors can be applied~\cite{ecss}.

\subsection*{Design assumptions (not yet fixed by the project)}

\begin{itemize}[nosep]
    \item Production hinge pin: \SI{12.7}{mm} (1/2\,in) diameter. The bench uses its measured 1/4\,in holes.
    \item Production harness routing: one gas-feed line to the inflatable tube at the root, plus one
    coolant or electrical jumper per bay, which is $n+1=8$ hinge-line crossings. The bench has the gas-feed
    line only.
    \item Deployed latch: one spring plunger (\SI{5}{N}, \SI{30}{\degree} ramp, \SI{25}{mm} from a centre
    pivot). One latch is enough because the scissor has a single degree of freedom.
\end{itemize}

% ====================================================
\section{Friction Source Table --- Production Array, On Orbit}
\label{sec:prod}

{\footnotesize\setlength{\tabcolsep}{3pt}
\renewcommand{\arraystretch}{1.2}
\begin{tabularx}{\linewidth}{@{} >{\raggedright\arraybackslash}p{2.6cm} c Y Y >{\raggedright\arraybackslash}p{1.8cm} r r r @{}}
\toprule
\textbf{Component} & \textbf{Class} & \textbf{Friction mechanism} &
\textbf{Coefficient / value (source)} & \textbf{Load basis} &
\multicolumn{3}{c}{\textbf{$Q$ [N\,m]}} \\
\cmidrule(l){6-8}
& & & & & Low & Nom. & High \\
\midrule
\makecell[tl]{Hinge pin joints\\($3n=21$)} & FR &
Journal sliding of pin in link bore; eq.~\eqref{eq:pins} &
$\mu$ = 0.01--0.04 sputtered MoS$_2$ in vacuum~\cite{spalvins};
0.61 bare Al--mild steel static~\cite{etb-friction} &
$N=2k\theta_0/L=\SI{0.356}{N}$ &
0.000904 & 0.00361 & 0.0551 \\
\makecell[tl]{Root guide /\\carriage (1)} & FR &
Sliding or rolling of the root top node along its guide; $Q=\mu N L$ &
$\mu$ = 0.002--0.003 ball-type LM guide~\cite{thk};
0.47 plain Al--steel slot, kinetic~\cite{etb-friction} &
same as pins &
0.00249 & 0.00374 & 0.586 \\
\makecell[tl]{Torsion springs\\($N_s=7$)} & HY &
Coil-to-coil and coil-to-arbor rubbing (hysteresis) &
5.9\% of spring torque, from the 10.2$\to$10.8 friction correction~\cite{engedge};
low/high $=0.5\times/2\times$. Coil spacing reduces it~\cite{leespring} &
$N_s k\theta_0=\SI{4.36}{N.m}$ &
0.128 & 0.257 & 0.513 \\
\makecell[tl]{Harness / flex\\lines (8 crossings)} & HA &
Bending and set of wires, hoses and gas line across each hinge line &
0.5 in-lb predicted, 3.2 in-lb measured, 20 in-lb worst ``zero-neutral'' per hinge line
(S-NPP array)~\cite{helfrich} &
per crossing &
0.452 & 2.89 & 18.1 \\
\makecell[tl]{Vacuum adhesion\\(breakaway only)} & HD &
Cold-welding adhesion at metal contacts after stowage in vacuum &
0 hard-anodised Al7075/SS; \SI{336}{mN} Alodine Al7075/Al7075;
\SI{1744}{mN} bare Al7075/Al7075~\cite{stm279} &
per contact, at $r_{\text{pin}}$ &
0 & 0.0853 & 0.443 \\
\makecell[tl]{Deployed latch\\(end of travel only)} & FR &
Plunger riding up the latch ramp;
$F_t = F_p\frac{\tan\alpha+\mu}{1-\mu\tan\alpha}$ &
$\mu$ = 0.42 steel--steel kinetic, 0.8 static~\cite{etb-friction};
latch geometry assumed &
$F_p=\SI{5}{N}$, $r=\SI{25}{mm}$ &
0.220 & 0.329 & 0.640 \\
\makecell[tl]{Inflatable kink tube} & -- &
Membrane rubbing on links, fold hysteresis &
Not quantified separately. Already inside the measured net $a_{\text{kink}}$ (Sec.~\ref{sec:unq}) &
-- & -- & -- & -- \\
\makecell[tl]{Hold-down / release} & -- &
Pin or bolt release friction &
Acts \emph{before} motion; sized with the release device's own margin &
-- & -- & -- & -- \\
\bottomrule
\end{tabularx}
}

\vspace{0.6em}
\textbf{Reading the table.} On orbit the harness and flex lines make up about 90\% of the nominal budget.
Pin-joint friction is almost negligible, because weightless pin loads are under \SI{1}{N} and the
lubricated $\mu$ is small. That changes only if the joints are bare metal or bind (see the thermal note in
Sec.~\ref{sec:unq}).

% ====================================================
\section{Friction Source Table --- Bench Test Article, 1\,g}
\label{sec:bench}

Same method with \texttt{testingParams.py}: $n=2$, $L=\SI{10}{in}$, 1/4\,in pins,
$M_j=\SI{1.062}{kg}$, $k=\SI{0.0248}{N.m/rad}$. Vacuum adhesion does not apply in air.

{\footnotesize\setlength{\tabcolsep}{3pt}
\renewcommand{\arraystretch}{1.2}
\begin{tabularx}{\linewidth}{@{} >{\raggedright\arraybackslash}p{2.6cm} c Y Y >{\raggedright\arraybackslash}p{1.8cm} r r r @{}}
\toprule
\textbf{Component} & \textbf{Class} & \textbf{Friction mechanism} &
\textbf{Coefficient / value (source)} & \textbf{Load basis} &
\multicolumn{3}{c}{\textbf{$Q$ [N\,m]}} \\
\cmidrule(l){6-8}
& & & & & Low & Nom. & High \\
\midrule
\makecell[tl]{Hinge pin joints\\($3n=6$)} & FR &
Journal sliding; eq.~\eqref{eq:pins}, multiplier 10 &
$\mu$ = 0.06--0.18 iglide J dry vs steel~\cite{igus};
0.61 bare Al--mild steel static~\cite{etb-friction} &
$N=M_j g=\SI{10.4}{N}$ &
0.0198 & 0.0595 & 0.202 \\
\makecell[tl]{Root guide /\\carriage (1)} & FR &
Guide at the root top node &
$\mu$ = 0.002--0.003 LM guide~\cite{thk}; 0.47 plain slot~\cite{etb-friction} &
$N=M_j g$ &
0.00529 & 0.00793 & 1.24 \\
\makecell[tl]{Torsion springs\\($N_s=2$)} & HY &
Coil / arbor rubbing &
5.9\%~\cite{engedge}; $0.5\times$/$2\times$ &
$N_s k\theta_0=\SI{0.078}{N.m}$ &
0.00229 & 0.00458 & 0.00917 \\
\makecell[tl]{Gas-feed line\\(1 crossing, if routed\\across a joint)} & HA &
Hose bending across the root hinge &
S-NPP per-hinge-line values~\cite{helfrich}; likely high for a small bench tube &
per crossing &
0.0565 & 0.362 & 2.26 \\
\makecell[tl]{Deployed latch\\(end of travel only)} & FR &
Plunger on ramp &
As production~\cite{etb-friction} &
$F_p=\SI{5}{N}$ &
0.220 & 0.329 & 0.640 \\
\bottomrule
\end{tabularx}
}

\vspace{0.6em}
\textbf{Reading the table.} In 1\,g the pin joints carry the panel weight, so pin friction is about
$16\times$ larger than on orbit. Any line crossing a bench joint, and the latch, would outweigh everything
else. Both should be measured directly on the bench rather than taken from flight literature.

% ====================================================
\section{System Friction Estimate}
\label{sec:sys}

Adhesion acts only at breakaway ($\theta_0$) and the latch acts only at the end of travel
($\theta_{\min}$), so the totals are split by where in the stroke they apply. The ECSS column applies
Table~4-2 of ECSS-E-ST-33-01C Rev.\,2: theoretical factor 3 on FR, HY, HA and HD, then the overall factor
of 2 in its minimum-actuation equation~\cite{ecss}:
\begin{equation}
T_{\min} = 2\,\big(3\,\text{FR} + 3\,\text{HY} + 3\,\text{HA} + 3\,\text{HD}\big).
\end{equation}
NASA-STD-5017B applies the same factor, $\text{FS}_v = 3.00$ on variable (friction and harness) torques from
analysis, in its torque-margin equation 4-1~\cite{nasa5017}.

The equivalent model inputs are $\tau_f = Q/N_s$ and $F_{\text{friction}} = \tau_f/r_o$, in the same
convention as \texttt{params.py}. The tip force is the same $Q$ seen as an axial push at the boom tip at
$\theta_1=\SI{75}{\degree}$: $F=Q/(nL\sin\theta_1)$.

\begin{table}[H]
\centering
\small
\renewcommand{\arraystretch}{1.2}
\begin{tabular}{@{} l l r r r r r r @{}}
\toprule
& & \multicolumn{3}{c}{\textbf{Raw $Q$ [N\,m]}} & \multicolumn{3}{c}{\textbf{ECSS-factored $Q$ [N\,m]}} \\
\cmidrule(lr){3-5}\cmidrule(l){6-8}
\textbf{Case} & \textbf{Stroke point} & Low & Nom. & High & Low & Nom. & High \\
\midrule
\multirow{3}{*}{\makecell[tl]{Production,\\orbit}}
 & Start ($\theta_0$) & 0.584 & 3.24 & 19.7 & 3.50 & 19.5 & 118 \\
 & Travel             & 0.584 & 3.16 & 19.2 & 3.50 & 18.9 & 115 \\
 & End ($\theta_{\min}$) & 0.803 & 3.49 & 19.9 & 4.82 & 20.9 & 119 \\
\midrule
\multirow{2}{*}{\makecell[tl]{Production, no lines\\across hinges}}
 & Travel & 0.132 & 0.264 & 1.15 & 0.791 & 1.58 & 6.92 \\
 & & & & & & & \\
\midrule
\multirow{2}{*}{\makecell[tl]{Bench, 1\,g}}
 & Travel, with gas line & 0.0839 & 0.434 & 3.71 & 0.503 & 2.60 & 22.3 \\
 & Travel, no line       & 0.0274 & 0.0720 & 1.45 & 0.165 & 0.432 & 8.72 \\
\bottomrule
\end{tabular}
\caption{System resistive torque. ``No lines'' rows drop the harness/flex-line term, which is the one
architecture choice that controls the total.}
\end{table}

\begin{table}[H]
\centering
\small
\renewcommand{\arraystretch}{1.2}
\begin{tabular}{@{} l r r r r l @{}}
\toprule
& & \multicolumn{3}{c}{\textbf{$F_{\text{friction}}$ [N]}} & \\
\cmidrule(lr){3-5}
\textbf{Case (travel)} & \textbf{$\tau_f$ nom. [N\,m]} & Low & Nom. & High & \textbf{Current model} \\
\midrule
Production, with lines & 0.451 & 1.64 & 8.88 & 54.1 & \SI{5}{N} (\SI{7.5}{N} with $\times1.5$) \\
Production, no lines   & 0.0377 & 0.37 & 0.74 & 3.23 & \SI{5}{N} (\SI{7.5}{N}) \\
Bench, with gas line   & 0.217 & 13.2 & 68.3 & 585 & \SI{5}{N} (\SI{7.5}{N}) \\
Bench, no line         & 0.0360 & 4.32 & 11.3 & 229 & \SI{5}{N} (\SI{7.5}{N}) \\
\midrule
\multicolumn{6}{@{}p{0.9\linewidth}}{Tip force at $\theta_1$ (low / nom. / high): production with lines
0.025 / 0.13 / 0.81\,N; bench 0.17 / 0.88 / 7.6\,N with line, 0.06 / 0.16 / 3.2\,N without.} \\
\bottomrule
\end{tabular}
\caption{Friction expressed in the \texttt{params.py} convention ($F_{\text{friction}}=\tau_f/r_o$). Note
$r_o=\SI{50.8}{mm}$ for production and \SI{3.175}{mm} for the bench, so the bench $F$ values are large only
because its moment arm is small.}
\end{table}

\subsection*{Final estimate}

\begin{center}
\fbox{\parbox{0.92\linewidth}{
\textbf{Production, on orbit:} nominal $Q \approx \SI{3.2}{N.m}$ (range 0.6--20), which is
$F_{\text{friction}} \approx \SI{9}{N}$ in the model's convention. About 90\% of that is harness and
flex-line drag. With no lines crossing the hinges, nominal $Q \approx \SI{0.26}{N.m}$
($F_{\text{friction}} \approx \SI{0.7}{N}$).\\[4pt]
\textbf{Bench, 1\,g, no line across the joints:} nominal $Q \approx \SI{0.072}{N.m}$
($F_{\text{friction}} \approx \SI{11}{N}$), rising to about \SI{0.40}{N.m} at the latch.}}
\end{center}

\subsection*{Check against the current spring}

\begin{itemize}[nosep]
    \item \textbf{Production.} At the current $k$, all springs together provide \SI{4.36}{N.m} at stow
    and \SI{3.64}{N.m} at $\theta_1$. The model budgets $N_s\tau_{f,\text{design}}=\SI{2.67}{N.m}$. The
    nominal estimate with lines is \SI{3.16}{N.m}, above the model's budget and only 15\% below the
    available spring torque at $\theta_1$. The spring still clears the \emph{raw} nominal friction. It does
    \textbf{not} clear the ECSS-factored requirement of about \SI{19.5}{N.m}, a shortfall of about
    $4.5\times$. If the harness and friction are measured, the factor drops from 3 to 1.5~\cite{ecss} and the
    requirement falls to about \SI{9.7}{N.m}, still $2.2\times$ the current spring.
    \item \textbf{Bench.} The spring from the \SI{5}{N} assumption gives \SI{0.065}{N.m} at $\theta_1$. The
    nominal no-line friction is \SI{0.072}{N.m}, so with the conservative 1\,g pin loads the bench article
    is at or below zero margin. The \SI{5}{N} (\SI{7.5}{N} factored) placeholder is too low for the bench;
    about \SI{11}{N} nominal is supported.
\end{itemize}

\subsection*{Self-consistent spring sizing}
\label{sec:sc}
Pin, guide and spring losses scale with $k$. Resizing the spring with the model's own rule
($\times1.5$ on the nominal start budget) and iterating to convergence gives $k=\SI{0.781}{N.m/rad}$ for
production. The nominal equivalent then settles at $F_{\text{friction}} \approx \SI{9.8}{N}$ (start), close to
the \SI{9.1}{N} above, because the harness term does not depend on $k$.

\section{Items Not Summed}
\label{sec:unq}

\begin{itemize}
    \item \textbf{Inflatable kink tube.} Membrane-on-link rubbing and fold hysteresis are real resistances.
    However, the Phase 2 model uses a \emph{net} kink moment $M_k=a\,\Delta P R^3$ with $a$ measured on the
    bench. Anything the tube loses to friction is already inside that measured $a$, so adding it here would
    count it twice. The single-kink bench test should therefore run the tube against representative link
    material.
    \item \textbf{Thermal clearance loss.} Aluminium (\SI{23e-6}{/K}) shrinks faster than carbon steel
    (10.8--\SI{12.5e-6}{/K})~\cite{etb-cte}. A steel pin in an Al bore therefore loses clearance when cold:
    \SI{7.7}{\micro m} on the 1/4\,in bench pin and \SI{15.5}{\micro m} on a 1/2\,in production pin for a
    \SI{100}{K} drop. Guzik and Benafan found the cold case, with a \SI{0.0076}{mm} pin gap, to be critical
    for their hinge~\cite{guzik}. This is a binding risk rather than a friction value. Set the minimum
    room-temperature clearance above the cold-case loss.
    \item \textbf{Hold-down and release.} Pin-puller or bolt friction acts before deployment starts and
    belongs to the release device's own force margin~\cite{guzik}.
    \item \textbf{Inertia.} This is not friction, and it is already shown to be negligible by the
    quasi-static check in the derivation.
\end{itemize}

\section{Recommendations}

\begin{enumerate}[nosep]
    \item Replace $F_{\text{friction}}=\SI{5}{N}$ in \texttt{params.py} with about \SI{9}{N} (production,
    lines crossing each bay) or about \SI{11}{N} (bench). Keep the $\times1.5$ factor until test data exist.
    \item Harness and flex-line drag decides the production budget. Route lines along the tube or through a
    single root crossing where possible, and measure the torque per crossing early. Measuring it also lets
    ECSS drop its factor from 3 to 1.5.
    \item Use a self-lubricating or MoS$_2$-coated pin interface with dissimilar or hard-anodised surfaces.
    Bare Al7075 against Al7075 gives the highest vacuum adhesion measured in~\cite{stm279}, and hard
    anodising prevents it.
    \item Use a ball-type guide at the root node. A plain slot moves the guide from about 0.1\% to 5\% of the
    production budget, and to about half of the bench budget in the high case.
    \item Measure the latch engagement torque on hardware. Guzik and Benafan report that analytical latch
    friction estimates did not match test data~\cite{guzik}.
\end{enumerate}

% ====================================================
\begin{thebibliography}{99}
\small

\bibitem{ecss}
ECSS Secretariat, \emph{ECSS-E-ST-33-01C Rev.\,2: Space Engineering --- Mechanisms}, ESA-ESTEC,
1~March~2019. Clause 4.7.5.3.1 and Table~4-2.
\url{https://ecss.nl/wp-content/uploads/2019/05/ECSS-E-ST-33-01C-Rev.2(1March2019).pdf}

\bibitem{nasa5017}
NASA, \emph{NASA-STD-5017B: Design and Development Requirements for Mechanisms}, approved 2022-12-06.
Section 4.3, Equation 4-1 and Table~1.
\url{https://standards.nasa.gov/sites/default/files/standards/NASA/B/2022-12-06-NASA-STD-5017B-Approved.pdf}

\bibitem{etb-friction}
The Engineering ToolBox, ``Friction and Friction Coefficients for Various Materials.''
\url{https://www.engineeringtoolbox.com/friction-coefficients-d_778.html}

\bibitem{igus}
igus GmbH, ``iglide J plastic bushings: Material data.''
\url{https://www.igus.com/plastic-bearings/resources/iglide-material-j}

\bibitem{spalvins}
T.~Spalvins, ``Lubrication with Sputtered MoS$_2$ Films: Principles, Operation, Limitations,''
NASA TM-105292, Lewis Research Center, 1991.
\url{https://ntrs.nasa.gov/citations/19910022902}

\bibitem{thk}
THK Co., Ltd., ``General Description --- Friction Coefficient,'' General Catalog, p.~B0-15, Table~4.
\url{https://tech.thk.com/es/products/pdf/en_b00_015.pdf}

\bibitem{engedge}
Engineers Edge, ``Torsion Spring Coil Deflection and Spring Rate Formulas.''
\url{https://www.engineersedge.com/analysis/torsion_spring_coil_deflection_16436.htm}

\bibitem{leespring}
Lee Spring, ``Learn About Torsion Springs.''
\url{https://www.leespring.com/learn-about-torsion-springs}

\bibitem{helfrich}
D.~Helfrich and A.~Sexton, ``Lessons Learned in the Flight Qualification of the S-NPP and NOAA-20 Solar
Array Mechanisms,'' \emph{Proc.\ 44th Aerospace Mechanisms Symposium}, NASA GRC, 2018.
\url{https://ntrs.nasa.gov/citations/20180003006}

\bibitem{stm279}
A.~Merstallinger, M.~Sales, E.~Semerad and B.\,D.~Dunn, \emph{Assessment of Cold Welding between
Separable Contact Surfaces due to Impact and Fretting under Vacuum}, ESA STM-279, November 2009.
\url{https://esamultimedia.esa.int/multimedia/publications/STM-279/STM-279.pdf}

\bibitem{guzik}
A.\,T.~Guzik and O.~Benafan, ``Design and Development of CubeSat Solar Array Deployment Mechanisms Using
Shape Memory Alloys,'' NASA/TM-2018-219914, Glenn Research Center, June 2018.
\url{https://ntrs.nasa.gov/citations/20180004164}

\bibitem{etb-cte}
The Engineering ToolBox, ``Thermal Expansion --- Linear Expansion Coefficients.''
\url{https://www.engineeringtoolbox.com/linear-expansion-coefficients-d_95.html}

\end{thebibliography}

\noindent\footnotesize All links were checked (HTTP 200) and the quoted values were read from each source on
29~September~2026.

\end{document}
